\documentclass[11pt,a4paper]{article}

% ======================================
% PAQUETES
% ======================================

\usepackage[utf8]{inputenc}
\usepackage[T1]{fontenc}
\usepackage[spanish,es-noshorthands]{babel}
\usepackage{lmodern}
\usepackage{microtype}
\usepackage{geometry}
\usepackage{graphicx}
\usepackage{booktabs}
\usepackage{tabularx}
\usepackage{array}
\usepackage{amsmath,amssymb}
\usepackage{float}
\usepackage{xcolor}
\usepackage{fancyhdr}
\usepackage{setspace}
\usepackage{caption}
\usepackage{hyperref}

\geometry{margin=2.5cm}
\graphicspath{{figuras/}}

\definecolor{academicblue}{RGB}{23,59,99}
\definecolor{academicgray}{RGB}{75,85,99}

\hypersetup{
    colorlinks=true,
    linkcolor=academicblue,
    citecolor=academicblue,
    urlcolor=academicblue,
    pdftitle={Valencia Airbnb Intelligence Lab},
    pdfauthor={Mat. Hector Roman}
}

\pagestyle{fancy}
\fancyhf{}
\fancyhead[L]{Valencia Airbnb Intelligence Lab}
\fancyhead[R]{Proyecto 1}
\fancyfoot[C]{\thepage}

\setlength{\headheight}{15pt}
\setlength{\parindent}{0pt}
\setlength{\parskip}{5pt}

\setstretch{1.25}

\captionsetup{
    font=small,
    labelfont=bf
}

\newcommand{\figura}[3]{
    \begin{figure}[H]
        \centering
        \includegraphics[width=#2\textwidth,height=0.68\textheight,keepaspectratio]{#1}
        \caption{#3}
    \end{figure}
}

% ======================================
% DOCUMENTO
% ======================================

\widowpenalty=10000
\clubpenalty=10000
\displaywidowpenalty=10000

\begin{document}

% ======================================
% PORTADA
% ======================================

\hypersetup{pageanchor=false}
\begin{titlepage}

\centering

{\Large\bfseries\color{academicblue}
UNIVERSIDAD YACHAY TECH\par}

\vspace{0.5cm}

{\large Maestría en Ciencia de Datos\par}

\vspace{0.3cm}

{\large PB -- Aprendizaje de Máquina\par}

\vfill

{\color{academicblue}\rule{\textwidth}{1.5pt}}

\vspace{0.7cm}

{\Huge\bfseries\color{academicblue}
Valencia Airbnb\\[0.3cm]
Intelligence Lab\par}

\vspace{0.6cm}

{\Large
Análisis estadístico, geoespacial y temporal
de alojamientos turísticos mediante Python\par}

\vspace{0.7cm}

{\color{academicblue}\rule{\textwidth}{1.5pt}}

\vfill

{\large
\begin{tabular}{rl}
\textbf{Autor:} & Mat. Héctor Román\\[0.3cm]
\textbf{Docente:} & Luis Fernando Zhinin Vera\\[0.3cm]
\textbf{Proyecto:} & 1\\[0.3cm]
\textbf{Versión:} & 2.1.0\\[0.3cm]
\textbf{Fecha:} & 26 de noviembre de 2026
\end{tabular}
}

\vfill

{\large
Informe científico-técnico\par}

\vspace{0.5cm}

\href{
https://github.com/alexroma1126-design/valencia-airbnb-intelligence-lab
}{Repositorio del proyecto en GitHub}

\end{titlepage}
\hypersetup{pageanchor=true}

% ======================================
% RESUMEN
% ======================================

\begingroup
% RESUMEN COMPACTO V21
\setstretch{1.10}
\small
\setlength{\parskip}{4pt}
\section*{Resumen}
\addcontentsline{toc}{section}{Resumen}

Este trabajo presenta Valencia Airbnb Intelligence Lab,
una plataforma computacional orientada a la exploración,
validación, análisis estadístico, análisis geoespacial
y estudio temporal de alojamientos turísticos en
Valencia, España.

La investigación utiliza datos procedentes de Inside
Airbnb, integrando tres conjuntos principales:
Listings, Calendar y Reviews, además de información
geográfica complementaria.

El sistema procesa 7\,968 alojamientos, 2\,908\,320
registros de disponibilidad y 509\,513 asociaciones
entre alojamientos y reseñas, distribuidos en 84 barrios.

La metodología comprende procesos de adquisición de
datos, auditoría de calidad, validación relacional,
estadística descriptiva, análisis de distribuciones,
medias móviles, descomposición STL y funciones de
autocorrelación.

Los resultados muestran una disponibilidad anunciada
general del 52,27\,\text{\char37} y un precio mediano registrado
de 134,10, expresado en las unidades monetarias de
la fuente.

Se identificaron 458\,154 identificadores únicos de
reseñas y asociaciones múltiples de determinados
identificadores con distintos alojamientos.

La plataforma incorpora ocho módulos interactivos
implementados mediante Streamlit y una arquitectura
modular en Python. La auditoría computacional superó
56 de 56 pruebas diseñadas, complementadas con
pruebas automatizadas de ejecución de la aplicación.

El trabajo establece una infraestructura reproducible
para investigaciones posteriores en aprendizaje de
máquina, sin atribuir capacidad predictiva a los
procedimientos exclusivamente exploratorios.

\textbf{Palabras clave:} ciencia de datos, Airbnb,
series temporales, STL, estadística descriptiva,
análisis geoespacial, Python.

% ======================================
% ABSTRACT
% ======================================

\section*{Abstract}
\addcontentsline{toc}{section}{Abstract}

This work presents Valencia Airbnb Intelligence Lab,
a computational platform for exploratory, statistical,
geospatial and temporal analysis of tourist accommodation
in Valencia, Spain.

The system integrates public datasets provided by Inside
Airbnb, comprising 7,968 accommodation listings,
2,908,320 calendar observations and 509,513 listing-review
associations across 84 represented neighbourhoods.

The methodology combines data acquisition, quality
assessment, relational validation, descriptive statistics,
moving averages, STL decomposition and autocorrelation
analysis.

The platform contains eight interactive Streamlit modules.
A computational audit successfully passed all 56 designed
validation tests, supplemented by automated application
execution tests.

The resulting infrastructure provides a reproducible
foundation for subsequent machine learning research.
No trained predictive model is claimed in the current
version.

\textbf{Keywords:} data science, Airbnb, time series,
geospatial analysis, Python, STL.

\endgroup

\clearpage

\tableofcontents

\clearpage

\listoffigures

\clearpage

% ======================================
% INTRODUCCION
% ======================================

\section{Introducción}

Las plataformas digitales de alojamiento generan
grandes cantidades de información relacionadas con
precios, disponibilidad, características estructurales
y comportamiento histórico de los usuarios.

El análisis computacional de estas fuentes permite
identificar patrones espaciales y temporales, comparar
zonas urbanas y construir indicadores que faciliten
la comprensión de la oferta turística.

Desde la perspectiva de la ciencia de datos,
la utilidad de un conjunto de información depende
no solamente de su volumen, sino también de su calidad,
consistencia, representatividad y estructura matemática.

Por esta razón, el presente proyecto desarrolla una
infraestructura analítica antes de incorporar algoritmos
predictivos.

El área de estudio corresponde a Valencia, considerando
la totalidad de los barrios representados en la
recopilación de datos.

El proyecto no se limita a una zona particular:
implementa mecanismos dinámicos de selección y
comparación entre barrios.

\subsection{Pregunta de investigación}

¿Cómo puede integrarse la información estructural,
geográfica y temporal de los alojamientos turísticos
de Valencia en una plataforma computacional
consistente, interpretable y reproducible?

\subsection{Justificación}

El estudio permite establecer una base metodológica
para investigaciones relacionadas con demanda turística,
distribución territorial de alojamientos, disponibilidad
anunciada y comportamiento temporal de las reseñas.

Asimismo, proporciona una infraestructura para futuros
problemas de clasificación, regresión, agrupamiento
y pronóstico temporal.

% ======================================
% OBJETIVOS
% ======================================

\section{Objetivos}

\subsection{Objetivo general}

Desarrollar una plataforma computacional reproducible
para el análisis exploratorio, estadístico, geoespacial
y temporal de alojamientos turísticos en Valencia,
utilizando Python y fundamentos matemáticos de
ciencia de datos.

\subsection{Objetivos específicos}

\begin{enumerate}
    \item Integrar los conjuntos Listings, Calendar
    y Reviews mediante identificadores consistentes.

    \item Evaluar la calidad, integridad y estructura
    de las bases de datos.

    \item Caracterizar estadísticamente los precios
    y tipos de alojamiento.

    \item Analizar la distribución espacial de la oferta
    turística.

    \item Construir series temporales de disponibilidad
    y reseñas registradas.

    \item Aplicar medias móviles, descomposición STL
    y análisis de autocorrelación.

    \item Implementar una aplicación interactiva
    para consultas y comparaciones entre barrios.

    \item Validar computacionalmente los procedimientos
    implementados y documentar su reproducibilidad.
\end{enumerate}

% ======================================
% DATOS
% ======================================

\section{Descripción de las fuentes de datos}

La fuente principal corresponde a Inside Airbnb
\cite{insideairbnb}.

Se utilizaron tres archivos detallados y dos fuentes
geográficas complementarias.

\begin{table}[H]
\centering
\caption{Dimensiones originales de las bases de datos.}
\begin{tabular}{lrr}
\toprule
Conjunto & Registros & Variables\\
\midrule
Listings & 7\,968 & 90\\
Calendar & 2\,908\,320 & 5\\
Reviews & 509\,513 & 6\\
\midrule
Total & 3\,425\,801 & --\\
\bottomrule
\end{tabular}
\end{table}

\subsection{Listings}

Cada registro corresponde a un anuncio de alojamiento
y contiene información estructural, geográfica,
económica y descriptiva.

La estructura puede representarse como

\[
L=(x_{ij})_{\substack{
1\leq i\leq N\\
1\leq j\leq p}},
\]

donde $N=7968$ y $p=90$.

Debido a la coexistencia de variables numéricas,
categóricas y textuales, la estructura no constituye
una matriz exclusivamente real.

\subsection{Calendar}

El calendario contiene información correspondiente
a 365 fechas para cada alojamiento.

Por tanto,

\[
7968\cdot365=2\,908\,320.
\]

El intervalo registrado se extiende desde el
25 de septiembre de 2026 hasta el
24 de septiembre de 2027.

Calendar no contiene precios diarios en la versión
analizada.

\subsection{Reviews}

La base contiene asociaciones entre alojamientos
y reseñas, con fechas comprendidas entre el
12 de octubre de 2010 y el 25 de septiembre de 2026.

Se identificaron:

\begin{itemize}
    \item 509\,513 asociaciones alojamiento-reseña.
    \item 458\,154 identificadores únicos de reseñas.
    \item 24\,699 identificadores asociados a más
    de un alojamiento.
\end{itemize}

Las ocurrencias adicionales corresponden a

\[
509513-458154=51359.
\]

Este resultado no implica necesariamente duplicación
errónea, puesto que los pares alojamiento-reseña
son únicos.

\subsection{Información geográfica}

Los archivos complementarios son
\texttt{neighbourhoods.csv} y
\texttt{neighbourhoods.geojson}.

El archivo GeoJSON contiene 88 elementos geográficos,
mientras que Listings presenta 84 barrios distintos.

Estas cantidades no deben interpretarse como
equivalentes sin comprobar la correspondencia
individual entre las entidades geográficas.

% ======================================
% METODOLOGIA
% ======================================

\section{Metodología computacional}

El proyecto utiliza una arquitectura modular
implementada en Python.

El procesamiento comprende las siguientes etapas:

\begin{enumerate}
    \item Adquisición automatizada de archivos.
    \item Verificación de integridad de archivos gzip.
    \item Carga y transformación de variables.
    \item Validación de identificadores y relaciones.
    \item Análisis exploratorio y estadístico.
    \item Construcción de indicadores geográficos
    y temporales.
    \item Implementación de módulos interactivos.
    \item Validación automatizada.
\end{enumerate}

\subsection{Integridad relacional}

Sean los conjuntos de identificadores

\[
I_L=\{i:i\in\mathrm{Listings}\},
\]

\[
I_C=\{i:i\in\mathrm{Calendar}\},
\]

\[
I_R=\{i:i\in\mathrm{Reviews}\}.
\]

Las pruebas permitieron comprobar que

\[
I_L=I_C,
\qquad
I_R\subseteq I_L.
\]

Además, se verificó la unicidad de las claves
compuestas

\[
(\mathrm{listing\_id},\mathrm{date})
\]

y

\[
(\mathrm{listing\_id},\mathrm{review\_id}).
\]

Esta propiedad resulta fundamental para evitar
multiplicaciones artificiales de registros durante
las operaciones relacionales.

\subsection{Calidad de datos}

Para cada variable $j$, se calcula la proporción
de observaciones faltantes mediante

\[
M_j=
\frac{100}{N}
\sum_{i=1}^{N}
\mathbf{1}_{\{x_{ij}\text{ faltante}\}}.
\]

\figura{01_missing_values.png}{0.95}{
Variables con mayor proporción de datos faltantes.
Elaboración propia a partir de Listings.}

La presencia de datos faltantes requiere decisiones
metodológicas dependientes del problema de investigación.

No se recomienda aplicar imputaciones o eliminaciones
generalizadas sin analizar el significado de cada variable.

\subsection{Procesamiento de precios}

Los precios originales se transformaron desde
representaciones textuales hacia valores numéricos.

Se conservaron los valores faltantes como observaciones
no disponibles, evitando generar precios artificiales.

El proceso obtuvo 7\,197 precios válidos.

% ======================================
% ESTADISTICA
% ======================================

\section{Análisis estadístico descriptivo}

Sea $X=\{x_1,\ldots,x_n\}$ una muestra de valores
numéricos válidos.

La media aritmética se define como

\[
\bar{x}=\frac{1}{n}\sum_{i=1}^{n}x_i.
\]

La varianza muestral es

\[
s^2=\frac{1}{n-1}
\sum_{i=1}^{n}(x_i-\bar{x})^2,
\]

y su desviación estándar corresponde a

\[
s=\sqrt{s^2}.
\]

La mediana es el cuantil de orden $0,5$:

\[
\widetilde{x}=Q_{0{,}5}.
\]

El rango intercuartílico se define como

\[
RIC=Q_{0{,}75}-Q_{0{,}25}.
\]

\subsection{Distribución de precios}

\begin{table}[H]
\centering
\caption{Estadísticos descriptivos de los precios registrados.}
\begin{tabular}{lr}
\toprule
Estadístico & Valor aproximado\\
\midrule
Observaciones válidas & 7\,197\\
Media & 166,70\\
Mediana & 134,10\\
Desviación estándar & 281,60\\
Mínimo & 9,13\\
Primer cuartil & 98,33\\
Tercer cuartil & 183,00\\
Máximo & 9\,900,00\\
\bottomrule
\end{tabular}
\end{table}

La media supera a la mediana y existe una diferencia
considerable entre el máximo y el tercer cuartil.

Estos resultados son compatibles con una distribución
asimétrica hacia valores elevados.

La existencia de precios extremos no demuestra por sí
misma que se trate de registros incorrectos.

\figura{04_price_distribution.png}{0.95}{
Distribución de precios hasta el percentil 99.
El recorte es exclusivamente visual.}

La figura excluye visualmente las observaciones
superiores al percentil 99 para facilitar la lectura.

Los estadísticos globales conservan todos los
precios válidos.


\paragraph{Interpretación de la distribución.}

La concentración de observaciones en intervalos
de precios relativamente bajos, junto con una cola
derecha prolongada, evidencia la importancia de
utilizar estadísticos robustos.

La mediana y el rango intercuartílico complementan
a la media y a la desviación estándar, especialmente
cuando existen observaciones de magnitud extrema.

\subsection{Comparación entre barrios}

\figura{02_neighbourhoods.png}{0.95}{
Veinte barrios con mayor número de alojamientos.}

Los cinco barrios con mayor representación son:

\begin{table}[H]
\centering
\caption{Barrios con mayor número de anuncios.}
\begin{tabular}{lr}
\toprule
Barrio & Alojamientos\\
\midrule
CABANYAL-CANYAMELAR & 774\\
RUSSAFA & 551\\
AIORA & 347\\
EL MERCAT & 340\\
EL CARME & 300\\
\bottomrule
\end{tabular}
\end{table}

Estas cantidades representan la oferta observada
en la fuente y no necesariamente la totalidad de
los alojamientos turísticos existentes.

\figura{05_price_boxplots.png}{0.98}{
Comparación mediante diagramas de caja.
Los valores atípicos se omiten solamente
de la representación visual.}

Los diagramas de caja permiten comparar medianas,
dispersión y rangos intercuartílicos.

La comparación es descriptiva y no constituye
una prueba formal de diferencias poblacionales.

\subsection{Tipos de alojamiento}

\figura{03_room_types.png}{0.9}{
Distribución según el tipo de alojamiento.}

\begin{table}[H]
\centering
\caption{Composición de la oferta registrada.}
\begin{tabular}{lr}
\toprule
Tipo & Cantidad\\
\midrule
Alojamiento completo & 6\,257\\
Habitación privada & 1\,682\\
Habitación compartida & 16\\
Habitación de hotel & 13\\
\bottomrule
\end{tabular}
\end{table}

% ======================================
% GEOGRAFIA
% ======================================

\section{Análisis geoespacial}

Cada alojamiento se caracteriza mediante un par
ordenado de coordenadas geográficas:

\[
P_i=(\lambda_i,\varphi_i),
\]

donde $\lambda_i$ representa la longitud y
$\varphi_i$ la latitud.

La distribución espacial permite estudiar la
concentración territorial de los anuncios y
su relación con los límites geográficos.

\figura{06_geographic_map.png}{0.85}{
Distribución espacial de los alojamientos.
Límites procedentes del archivo GeoJSON.}

La plataforma utiliza Folium \cite{folium} y Leaflet para
la visualización interactiva, además de una
alternativa basada directamente en coordenadas.

El mapa estático del informe se construye
sin servicios cartográficos externos.

Las coordenadas proporcionadas por la fuente pueden
ser aproximadas, por lo que no deben interpretarse
como localizaciones exactas de inmuebles.

% ======================================
% DISPONIBILIDAD
% ======================================


\paragraph{Interpretación territorial.}

La representación geográfica muestra una mayor
concentración visual de anuncios en el núcleo
urbano de Valencia, con agrupaciones adicionales
en sectores próximos al litoral.

La intensidad visual de los puntos representa
la concentración de anuncios observados y no una
densidad normalizada por superficie.

Una estimación formal de densidad espacial
requeriría calcular áreas geográficas o incorporar
métodos específicos de estimación espacial.

\section{Modelización de la disponibilidad}

La información de Calendar permite construir
una matriz binaria

\[
A\in\{0,1\}^{N\times T},
\]

donde $N=7968$ y $T=365$.

Definimos

\[
A_{it}=
\begin{cases}
1,&\text{si el anuncio aparece disponible},\\
0,&\text{si aparece no disponible}.
\end{cases}
\]

La disponibilidad diaria se calcula como

\[
D_t=\frac{100}{N}\sum_{i=1}^{N}A_{it}.
\]

El indicador general corresponde a

\[
D=
\frac{100}{NT}
\sum_{i=1}^{N}\sum_{t=1}^{T}A_{it}.
\]

El resultado registrado fue

\[
D=52{,}27\;\text{\char37}.
\]

\figura{07_daily_availability.png}{0.98}{
Disponibilidad diaria anunciada y media móvil
de siete días.}

La disponibilidad máxima observada corresponde
a 5\,889 anuncios en una fecha, mientras que
la mínima corresponde a 528.

Estos indicadores describen la información
publicada en el calendario, no la ocupación real.

\subsection{Disponibilidad durante estancias completas}

Para una estancia de $k$ noches con fecha de
entrada $t$, definimos

\[
B_i(t,k)=
\prod_{j=0}^{k-1}A_{i,t+j}.
\]

Como cada elemento es binario,

\[
B_i(t,k)=1
\]

si y solamente si todas las noches del intervalo
aparecen disponibles.

Además, deben verificarse las restricciones de
estancia mínima y máxima de la fecha de entrada:

\[
m_{it}\leq k\leq M_{it}.
\]

El indicador de compatibilidad se representa como

\[
C_i(t,k)=
B_i(t,k)
\mathbf{1}_{\{m_{it}\leq k\leq M_{it}\}}.
\]

El número de anuncios compatibles es

\[
N_c(t,k)=\sum_{i=1}^{N}C_i(t,k).
\]

Para una estancia de tres noches desde el
14 de diciembre de 2026, se obtuvieron
3\,467 alojamientos compatibles.

Este resultado fue comprobado mediante dos
procedimientos computacionales independientes.

No constituye una garantía de disponibilidad
efectiva en tiempo real ni una confirmación
de reserva.

% ======================================
% SERIES TEMPORALES
% ======================================

\section{Análisis de series temporales}

Una serie temporal es una sucesión ordenada
de observaciones

\[
\{Y_t:t\in\mathcal{T}\},
\]

donde el índice representa el tiempo.

En este proyecto se estudian principalmente
las series de disponibilidad anunciada y
las reseñas registradas.

\subsection{Construcción de la serie de reseñas}

Sea $\mathcal{R}_t$ el conjunto de registros
asociados a un período temporal $t$.

El número de asociaciones es

\[
Y_t^{(a)}=|\mathcal{R}_t|.
\]

El número de identificadores únicos es

\[
Y_t^{(u)}=
\left|
\{r.\mathrm{id}:r\in\mathcal{R}_t\}
\right|.
\]

En consecuencia,

\[
Y_t^{(u)}\leq Y_t^{(a)}.
\]

La distinción evita interpretar automáticamente
cada asociación como una reseña independiente.

\figura{08_monthly_reviews.png}{0.98}{
Evolución mensual de asociaciones e identificadores
únicos de reseñas.}

Se excluyeron los meses inicial y final de manera
conservadora para evitar posibles períodos incompletos.

La serie resultante contiene 190 observaciones mensuales.


\paragraph{Comportamiento histórico observado.}

La evolución mensual evidencia un crecimiento
de las reseñas registradas durante gran parte
del intervalo analizado, una reducción pronunciada
alrededor de 2020 y una recuperación posterior.

En los últimos años también se observa un incremento
de la amplitud de las fluctuaciones temporales.

Estos resultados deben interpretarse considerando
que la serie histórica se construye a partir
de los anuncios incluidos en la recopilación
analizada.

Por tanto, la evolución observada no permite
estimar directamente el crecimiento histórico
de todas las reservas turísticas de Valencia.

\subsection{Medias móviles}

La media móvil retrospectiva de orden $k$ se
define como

\[
MA_t^{(k)}=
\frac{1}{k}\sum_{j=0}^{k-1}Y_{t-j}.
\]

Este operador reduce fluctuaciones locales
y facilita la identificación visual de tendencias.

Una ventana de doce meses puede ayudar a
suavizar variaciones estacionales anuales
en series de frecuencia mensual.

\subsection{Descomposición STL}

El método STL, desarrollado por Cleveland y
colaboradores \cite{cleveland1990}, permite
descomponer una serie temporal mediante
suavizado local.

La representación aditiva se expresa como

\[
Y_t=T_t+S_t+R_t,
\]

donde

\begin{itemize}
    \item $T_t$ representa la tendencia.
    \item $S_t$ representa la componente estacional.
    \item $R_t$ representa el residuo.
\end{itemize}

Para la serie mensual se empleó un período
estacional de doce observaciones.

\figura{09_stl_decomposition.png}{0.98}{
Descomposición STL robusta de la serie mensual
de identificadores únicos de reseñas.}

El período seleccionado constituye una hipótesis
exploratoria de estacionalidad anual.

La descomposición no implica que los patrones
se mantengan constantes en el futuro.


\paragraph{Interpretación de los componentes.}

La tendencia estimada mediante STL refleja
el crecimiento de largo plazo de las reseñas
registradas, incluyendo una reducción temporal
y una recuperación posterior.

La componente estacional presenta fluctuaciones
de amplitud creciente en los años recientes.

Este comportamiento motiva la exploración futura
de transformaciones que estabilicen la varianza,
sin que su conveniencia pueda establecerse
exclusivamente mediante inspección visual.

Los residuos representan las variaciones no
explicadas por las componentes estimadas de
tendencia y estacionalidad.

\subsection{Validación de la reconstrucción}

Se verificó el error máximo

\[
\varepsilon_{\max}=
\max_t
\left|
Y_t-(T_t+S_t+R_t)
\right|.
\]

El resultado de la auditoría fue

\[
\varepsilon_{\max}
=9{,}095\times10^{-13}.
\]

Este valor es compatible con el error numérico
de las operaciones en punto flotante.

La prueba confirma la consistencia de la
reconstrucción aditiva, pero no constituye
una evaluación de precisión predictiva.

% ======================================
% ACF Y PACF
% ======================================

\section{Autocorrelación y dependencia temporal}

\subsection{Función de autocorrelación}

La función de autocorrelación teórica de una
serie débilmente estacionaria se define como

\[
\rho(k)=
\frac{\operatorname{Cov}(Y_t,Y_{t-k})}
{\operatorname{Var}(Y_t)}.
\]

Un estimador muestral es

\[
\widehat{\rho}(k)=
\frac{
\sum_{t=k+1}^{n}
(Y_t-\bar{Y})(Y_{t-k}-\bar{Y})
}{
\sum_{t=1}^{n}
(Y_t-\bar{Y})^2
}.
\]

La autocorrelación permite estudiar la dependencia
lineal entre observaciones separadas por $k$
períodos temporales.

\subsection{Autocorrelación parcial}

La función de autocorrelación parcial, PACF,
mide la asociación lineal entre $Y_t$ y $Y_{t-k}$
después de controlar los retardos intermedios.

Puede formularse mediante una regresión
autorregresiva de orden $k$:

\[
Y_t=
c+\sum_{j=1}^{k}
\phi_{kj}Y_{t-j}
+\varepsilon_t.
\]

El último coeficiente $\phi_{kk}$ representa
la autocorrelación parcial bajo la formulación
autorregresiva correspondiente.

\figura{10_acf_pacf.png}{0.98}{
Funciones ACF y PACF de la serie mensual
de reseñas registradas.}

Estas funciones se incluyen como herramientas
de diagnóstico exploratorio.

La interpretación para identificar órdenes de
modelos ARIMA requiere evaluar previamente
estacionariedad, tendencia y estacionalidad.

Los gráficos de la serie original no constituyen
por sí solos evidencia suficiente para seleccionar
un modelo definitivo.

% ======================================
% APLICACION
% ======================================


\subsection{Interpretación de los resultados}

La función de autocorrelación de la serie mensual
presenta valores positivos elevados en los primeros
retardos y un decaimiento relativamente lento.

Por otra parte, la función de autocorrelación
parcial muestra una asociación especialmente
elevada en el primer retardo.

Este comportamiento es compatible con una fuerte
persistencia temporal y con la presencia de una
tendencia en la serie original.

Sin embargo, no constituye evidencia suficiente
para establecer un modelo autorregresivo de
orden uno ni para confirmar la estacionariedad.

Una etapa predictiva posterior debería incorporar
transformaciones, pruebas de estacionariedad,
diagnóstico de residuos y validación cronológica,
siguiendo una metodología de pronóstico
estadístico \cite{hyndman2021}.

Los procedimientos computacionales de diagnóstico
se implementaron mediante Statsmodels
\cite{statsmodels}.

\section{Desarrollo de la plataforma interactiva}

La plataforma se implementó mediante Streamlit
\cite{streamlit}, Pandas \cite{mckinney2010}
y herramientas complementarias de análisis
numérico y visualización.

\begin{table}[H]
\centering
\caption{Arquitectura funcional de la aplicación.}
\begin{tabularx}{\textwidth}{
p{0.05\textwidth}
p{0.31\textwidth}
X}
\toprule
N. & Módulo & Funcionalidad\\
\midrule
1 & Exploración &
Calidad de datos, dimensiones e integridad.\\
2 & Barrios y mapas &
Distribución territorial de anuncios.\\
3 & Precios &
Distribuciones y estadística descriptiva.\\
4 & Series temporales &
Disponibilidad e historial de reseñas.\\
5 & Consultas rápidas &
Filtros de características y fechas.\\
6 & Laboratorio matemático &
STL, medias móviles, ACF y PACF.\\
7 & Comparador de barrios &
Comparación de indicadores descriptivos.\\
8 & Estancias completas &
Disponibilidad consecutiva y restricciones.\\
\bottomrule
\end{tabularx}
\end{table}

\subsection{Arquitectura del código}

La estructura principal separa la lógica analítica
de la representación visual:

\begin{verbatim}
app/
    app.py
    advanced_tabs.py
    geographic.py

src/
    analytics.py
    download_data.py
    inspect_data.py
    query_engine.py
    audit_reviews.py

results/
    tables/

test_analytics.py
test_streamlit.py
validate_project.py
\end{verbatim}

Esta separación facilita el mantenimiento,
la extensión y la evaluación independiente
de los procedimientos matemáticos.

% ======================================
% VALIDACION
% ======================================

\section{Validación computacional}

La validación se diseñó para comprobar propiedades
estructurales, relaciones entre conjuntos de datos
y consistencia de los algoritmos.

\begin{table}[H]
\centering
\caption{Resultados de la auditoría integral.}
\begin{tabular}{lr}
\toprule
Indicador & Resultado\\
\midrule
Pruebas ejecutadas & 56\\
Pruebas superadas & 56\\
Pruebas fallidas & 0\\
Tiempo observado & 4,03 s\\
Porcentaje de pruebas superadas & 100\,\text{\char37}\\
\bottomrule
\end{tabular}
\end{table}

La validación incluyó:

\begin{itemize}
    \item Existencia de archivos requeridos.
    \item Integridad de archivos comprimidos.
    \item Dimensiones de las bases.
    \item Unicidad de identificadores.
    \item Correspondencia entre tablas.
    \item Consistencia de claves compuestas.
    \item Validación de rangos geográficos.
    \item Consultas por barrio.
    \item Funciones estadísticas.
    \item Reconstrucción STL.
    \item Propiedades de ACF y PACF en retardo cero.
    \item Disponibilidad durante estancias completas.
\end{itemize}

\subsection{Validación de la aplicación}

Las pruebas automatizadas de Streamlit comprobaron
la ejecución inicial, la presencia de ocho módulos
y la selección dinámica de cinco barrios:

\begin{center}
BENIFERRI, BENICALAP, RUSSAFA, CAMPANAR y ALBORS.
\end{center}

No se registraron excepciones en las ejecuciones
comprobadas.

La validación automatizada no garantiza la correcta
representación de todos los componentes JavaScript
ni de todas las combinaciones posibles de entradas.

\subsection{Integridad criptográfica}

Los archivos principales fueron registrados mediante
huellas digitales SHA-256.

Para un archivo $F$, se define

\[
h=\operatorname{SHA256}(F).
\]

La comparación posterior de estas huellas permite
detectar modificaciones de los archivos originales.

Los resultados se almacenan en
\texttt{data\_hashes.json}.

% ======================================
% APRENDIZAJE DE MAQUINA
% ======================================

\section{Relación con el aprendizaje de máquina}

El presente trabajo desarrolla una infraestructura
de datos y análisis necesaria para abordar
problemas de aprendizaje de máquina.

Formalmente, un problema de aprendizaje supervisado
puede expresarse mediante una muestra

\[
\mathcal{D}=
\{(\mathbf{x}_i,y_i)\}_{i=1}^{n},
\]

y un modelo parametrizado $f_{\theta}$.

El entrenamiento busca minimizar una función
de pérdida empírica:

\[
\widehat{\theta}
=
\operatorname*{arg\,min}_{\theta}
\frac{1}{n}
\sum_{i=1}^{n}
\mathcal{L}
\left(
y_i,f_{\theta}(\mathbf{x}_i)
\right).
\]

La plataforma desarrollada facilita la construcción
de los conjuntos de características requeridos
por este proceso.

Entre las aplicaciones futuras se encuentran:

\begin{itemize}
    \item Regresión para estimación de precios.
    \item Clasificación de alojamientos.
    \item Agrupamiento de anuncios o barrios.
    \item Pronóstico temporal de reseñas.
    \item Identificación de patrones espaciales.
\end{itemize}

\subsection{Posible extensión hacia modelos temporales}

Un modelo SARIMA puede representarse mediante

\[
\Phi_P(B^s)\phi_p(B)
(1-B)^d(1-B^s)^D Y_t
=
\Theta_Q(B^s)\theta_q(B)\varepsilon_t,
\]

donde $B$ representa el operador de retardo
y $s$ el período estacional.

Esta familia de modelos constituye una posible
extensión metodológica \cite{hyndman2021}.

Su incorporación requeriría evaluar estacionariedad,
seleccionar parámetros y comparar resultados
mediante validación cronológica.

No se entrenaron modelos SARIMA ni se estimaron
sus capacidades predictivas en la versión 2.0.0.

\subsection{Evaluación predictiva futura}

Para predicciones $\widehat{y}_i$ y observaciones
reales $y_i$, podrían utilizarse las métricas

\[
\mathrm{MAE}=
\frac{1}{n}\sum_{i=1}^{n}
|y_i-\widehat{y}_i|,
\]

\[
\mathrm{RMSE}=
\sqrt{
\frac{1}{n}
\sum_{i=1}^{n}
(y_i-\widehat{y}_i)^2
}.
\]

En problemas temporales, la partición de entrenamiento
y evaluación debe respetar el orden cronológico
para evitar utilizar información futura.

% ======================================
% LIMITACIONES
% ======================================

\section{Limitaciones metodológicas}

El análisis presenta las siguientes limitaciones:

\begin{enumerate}
    \item La disponibilidad publicada no equivale
    necesariamente a la ocupación real.

    \item Las reseñas no representan directamente
    el número total de huéspedes o reservas.

    \item Un identificador de reseña puede asociarse
    a diferentes alojamientos.

    \item Los precios registrados no representan
    necesariamente precios efectivos de contratación.

    \item La moneda del precio requiere confirmación
    independiente mediante documentación de la fuente.

    \item El calendario analizado no contiene
    precios diarios.

    \item Las coordenadas geográficas pueden
    ser aproximadas.

    \item El historial de reseñas está condicionado
    por los anuncios representados en la recopilación,
    lo que puede introducir sesgos de selección.

    \item La comparación descriptiva no implica
    causalidad entre características de los barrios
    y comportamiento turístico.

    \item Las pruebas automatizadas no sustituyen
    totalmente la validación visual ni la evaluación
    de generalización predictiva.
\end{enumerate}

% ======================================
% REPRODUCIBILIDAD
% ======================================

\section{Reproducibilidad científica}

El proyecto \cite{project} utiliza Git como sistema de control
de versiones y GitHub como repositorio remoto.

La versión analítica inicial fue registrada
mediante la etiqueta

\begin{center}
\texttt{v2.0.0}
\end{center}

y el identificador abreviado de commit

\begin{center}
\texttt{95a1e84}.
\end{center}

El repositorio contiene el código fuente,
la documentación, las dependencias registradas
y los resultados de las auditorías.

\subsection{Repositorio}

\url{https://github.com/alexroma1126-design/valencia-airbnb-intelligence-lab}

El repositorio fue configurado inicialmente como
privado, por lo que su consulta requiere autorización.

\subsection{Ejecución}

Las instrucciones básicas para Windows son:

\begin{verbatim}
python -m venv .venv

.\.venv\Scripts\Activate.ps1

python -m pip install -r requirements-lock.txt

python src/download_data.py

python validate_project.py

python test_streamlit.py

python -m streamlit run app/app.py
\end{verbatim}

La reproducción exacta de los resultados requiere
recuperar los mismos archivos de origen y comprobar
sus huellas SHA-256.

Se verificó la clonación del repositorio remoto
y la compilación sintáctica de sus módulos Python.

La ejecución completa desde un entorno independiente
y recién instalado permanece como una validación
adicional no realizada.

% ======================================
% CONCLUSIONES
% ======================================

\section{Conclusiones}

El desarrollo de Valencia Airbnb Intelligence Lab
permitió integrar 3\,425\,801 registros procedentes
de tres conjuntos de información complementarios.

La validación relacional confirmó la consistencia
entre los identificadores de alojamientos y las
observaciones correspondientes al calendario
y las reseñas.

El análisis exploratorio permitió caracterizar
la distribución de precios, la composición de
los tipos de alojamiento y la concentración
territorial de anuncios.

La mediana de precios registrada fue de 134,10,
mientras que la disponibilidad anunciada general
alcanzó el 52,27\,\text{\char37}.

La distinción entre asociaciones alojamiento-reseña
e identificadores únicos constituye un resultado
metodológico relevante para la construcción
de indicadores temporales.

La implementación de medias móviles, descomposición
STL y funciones de autocorrelación proporciona
herramientas para investigar la estructura temporal
de la información.

La auditoría computacional superó las 56 pruebas
diseñadas y la aplicación integró ocho módulos
interactivos.

Finalmente, el trabajo establece una plataforma
analítica extensible, documentada y versionada,
capaz de servir como fundamento para posteriores
investigaciones de aprendizaje de máquina.

% ======================================
% BIBLIOGRAFIA
% ======================================

\clearpage

\begin{thebibliography}{99}

\bibitem{insideairbnb}
Inside Airbnb.
\textit{Get the Data}.
\url{https://insideairbnb.com/get-the-data/}.

\bibitem{cleveland1990}
Cleveland, R. B., Cleveland, W. S.,
McRae, J. E. y Terpenning, I. (1990).
STL: A seasonal-trend decomposition procedure
based on loess.
\textit{Journal of Official Statistics},
6(1), 3--73.

\bibitem{hyndman2021}
Hyndman, R. J. y Athanasopoulos, G. (2021).
\textit{Forecasting: Principles and Practice},
3rd edition. OTexts.
\url{https://otexts.com/fpp3/}.

\bibitem{mckinney2010}
McKinney, W. (2010).
Data structures for statistical computing
in Python.
\textit{Proceedings of the 9th Python in
Science Conference}, 56--61.

\bibitem{statsmodels}
Statsmodels.
\textit{Time Series Analysis Documentation}.
\url{https://www.statsmodels.org/}.

\bibitem{streamlit}
Streamlit.
\textit{Official Documentation}.
\url{https://docs.streamlit.io/}.

\bibitem{folium}
Folium.
\textit{Documentation}.
\url{https://python-visualization.github.io/folium/}.

\bibitem{project}
Román, H. (2026).
\textit{Valencia Airbnb Intelligence Lab,
versión 2.0.0}.
Repositorio de código fuente.
\url{https://github.com/alexroma1126-design/valencia-airbnb-intelligence-lab}.

\end{thebibliography}

\end{document}
